\subsection{Bases in Hilbert Spaces}\label{sec:bases-hilbert-spaces}

\subsubsection{The Existence of a Basis}\label{subsec:hilbert-basis-existence}

From linear algebra, we know that every Euclidean vector space has a Cartesian basis. In this subsection, we shall extend this result to Hilbert spaces. We shall need the following definition.

\begin{Definition}
\label{def:unordered-hilbert-series-sum}
Let $H$ be a Hilbert space and $I$ a (possibly uncountable) index set. Let $\{x_i\}_{i\in I}$ be a family of elements of $H$. We say that $\sum_{i\in I}x_i=x$ if at most countably many of the $x_i$ are nonzero, and if for any enumeration of these nonzero elements we have $x=\sum_{j\in\mathbb{N}}x_{i(j)}$.
\end{Definition}

\begin{Remark}
\label{rem:unordered-series-countability}
Note that while it is convenient for us to allow for the possibility of an uncountable index set, at most countably many elements can be nonzero if this notion of convergence is to make sense. To see this, note that for any series of real numbers to be absolutely convergent, it can have at most a finite number of terms with norm greater than, say, $1/n$ for and natural number $n$. Hence, it can have at most countably many nonzero terms. The above definition of convergence is a generalization of absolute convergence of a series of real number, and the following conditions are easily shown to be equivalent to that definition. We leave the proof to the reader (Problem~\ref{ex:hilbert-series-convergence-criteria}).
\end{Remark}

\begin{Lemma}
\label{lem:unordered-hilbert-series-convergence-criteria}
The sum $x=\sum_{i\in I}x_i$ exists if and only if either of the following hold:\begin{enumerate}
\item For every $\epsilon>0$, there is a finite subset $J_\epsilon$ of $I$ such that for any finite $J$ with $J_\epsilon\subset J\subset I$ we have
\begin{equation}\label{eq:hilbert-series-finite-sum-approximation}
\left\|x-\sum_{i\in J}x_i\right\|<\epsilon.
\end{equation}

\item For every $\epsilon>0$ there is a finite subset $J_\epsilon$ of $I$ such that
\begin{equation}\label{eq:hilbert-series-finite-tail-cauchy-condition}
\left\|\sum_{i\in J}x_i\right\|<\epsilon
\end{equation}
for any finite subset $J$ of $I$ with $J\cap J_\epsilon=\emptyset$.\end{enumerate}
\end{Lemma}

In the following, we are interested in sums of orthogonal elements. We have the following lemma.

\begin{Lemma}
\label{lem:orthogonal-series-summability}
Let $\{x_i\}_{i\in I}$ be a family of mutually orthogonal elements of a Hilbert space $H$. Then $\sum_{i\in I}x_i$ exists if and only if $\sum_{i\in I}\|x_i\|^2<\infty$. In this case we have, moreover,
\begin{equation}\label{eq:orthogonal-series-pythagorean-norm}
\left\|\sum_{i\in I}x_i\right\|^2=\sum_{i\in I}\|x_i\|^2.
\end{equation}

\end{Lemma}

\begin{Proof}
For any finite subset $J$ of $I$ we use the fact that elements of $\{x_i\}_{i\in I}$ are mutually orthogonal to get
\begin{equation}\label{eq:finite-orthogonal-sum-norm-calculation}
\left\|\sum_{i\in J}x_i\right\|^2=\left(\sum_{i\in J}x_i,\sum_{l\in J}x_l\right)=\sum_{i\in J}(x_i,x_i)=\sum_{i\in J}\|x_i\|^2.
\end{equation}
The rest follows from Lemma~\ref{lem:unordered-hilbert-series-convergence-criteria}.
\end{Proof}

\begin{Definition}
\label{def:orthonormal-family}
A family $\{x_i\}_{i\in I}$ of mutually orthogonal elements of $H$ is called \textbf{orthonormal} if $\|x_i\|=1$ for every $i\in I$.
\end{Definition}
